\section{Preliminaries}
\subsection{$\Gamma$-calculus on Boolean hypercube}
Let $\Omega = \{\pm 1\}^n$ be the boolean hypercube.
For any $x \in \Omega$, we use $x^i$ to denote the vector obtained by flipping the $i$-th coordinate of $x$ and $x^{ij} := (x^i)^j$.
Given the transition rates $\{q_i: \Omega \mapsto \R_+\}_{i\in[n]}$, we define the operator $L$ on functions $f:\Omega \to \R$ as follows:
\begin{align*}
    \forall x \in \Omega,\quad Lf(x) = \sum_{i=1}^n q_i(x)(f(x^i) - f(x)).
\end{align*}
The \textit{continuous-time Markov semigroup} $(P_t)_{t\geq 0}$ is given by $P_t = \e^{t L}$.
Let $\mu$ be the \textit{stationary distribution} of the Markov semigroup $(P_t)_{t\geq 0}$.
And in this paper, we assume that $L$ is \textit{reversible} with respect to $\mu$, that is, for any $x\in \Omega$ and $i\in [n]$, it holds 
$$\mu(x)q_i(x) = \mu(x^i)q_i(x^i).$$

\begin{definition}[Weighted carr\'e du champ operator $\Gamma_w$]
    Given the weight functions $w_i: \Omega \to \R_+$ for every $i\in [n]$, the \textit{weighted carr\'e du champ operator} $\Gamma_w$ is defined that for any function $f,g:\Omega \to \R$,
    \begin{align*}
        \Gamma_w(f,g) := \frac 12 (L_w(fg)-gL_wf-fL_wg),
    \end{align*}
    where $L_w$ is the operator defined by the reweighted transition rates $\{w_iq_i\}_{i\in[n]}$ (for which $\mu$ is not necessarily stationary). 
    We write $\Gamma_w(f) := \Gamma_w(f,f)$ for simplicity.
    By a straightforward calculation, for any $x \in \Omega$, it holds
    \begin{align*}
       \Gamma_w(f,g)(x) = \frac 12\sum_{i=1}^n w_i(x)q_i(x)(f(x^i) - f(x))(g(x^i) - g(x)).
    \end{align*}
    We also use $\Gamma_{w,i}$ to denote the contribution of the $i$-th coordinate to $\Gamma_w$, that is, 
    \begin{align*}
        \forall x \in \Omega,\quad \Gamma_{w,i}(f,g)(x) := \frac 12 w_i(x)q_i(x)(f(x^i) - f(x))(g(x^i) - g(x)).
    \end{align*}
    We write $\Gamma_{w,i}(f) := \Gamma_{w,i}(f,f)$ for simplicity.
    Note, $\Gamma_w(f) = \sum_{i=1}^n \Gamma_{w,i}(f)$.
\end{definition}

We omit the subscript $w$ when $w_i \equiv 1$ for all $i \in [n]$.

It is also convenient to use the weighted discrete gradient to represent the carr\'e du champ operator.


\begin{definition}[Weighted discrete gradient of $f$]\label{def:weighted-gradient}
    The \textit{weighted discrete gradient} of function $f: \Omega \to \mathbb{R}$ is defined as the vector 
    \begin{align*}
        \nabla_w f := (\delta_{w,1} f, \ldots, \delta_{w,n} f),\quad \text{where} \quad \delta_{w,i} f := \sqrt{\frac 12w_i(x)q_i(x)}(f(x^i) - f(x)).
    \end{align*}
    The \textit{absolute weighted discrete gradient} is the entrywise absolute value of $\nabla_w f$, denoted by $|\nabla_w| f := (|\delta_{w,1}| f, \ldots, |\delta_{w,n}| f)$.
\end{definition}

\begin{fact}
    We have that $\Gamma_{w,i}(f,g) = (\delta_{w,i} f) (\delta_{w,i} g)$ and $\Gamma_w(f,g) = \inner{\nabla_w f}{\nabla_w g}$.
    In particular, $\Gamma_{w,i}(f) = (\delta_{w,i} f)^2$ and $\Gamma_w(f) = \norm{\nabla_w f}_2^2$.
\end{fact}

\ZC{$|\nabla_w f|$ vs $|\nabla_w| f$, $|\delta_{w,1} f|$ vs $|\delta_{w,1}| f$}

\ZJ{I used $|\nabla_w| f$ and $|\delta_{w,1}| f$ in Section 5.}

\ZC{Mention the definition of discrete gradient is slightly non-standard.}

\ZC{add citation for weighted $\Gamma$}

\begin{definition}[Second-order weighted carr\'e du champ operator $\Gamma_{2;w}$]
    The \textit{second-order weighted carr\'e du champ operator} $\Gamma_{2;w}$ satisfies that for any function $f:\Omega \to \R$,
    \begin{align*}
        \Gamma_{2;w}(f) := \frac 12 L\Gamma_w(f) - \Gamma_w(f,Lf).
    \end{align*}
\end{definition}

\begin{lemma}[\cite{BHLLMY15,FS18}, Chain rule]
\label{lem:chain-rule}
    For any function $f: \Omega \to \R$, it holds
    \begin{align*}
        \sqrt{f} L \sqrt{f} = \frac{1}{2} L f - \Gamma(\sqrt{f}).
    \end{align*}
\end{lemma}

% \begin{proof}
%     For any $x \in \Omega$, by definition, we have
%     \begin{align*}
%         \sqrt{f(x)} L \sqrt{f}(x) &= \sum_{i=1}^n q_i(x) \sqrt{f(x)} (\sqrt{f(x^i)} - \sqrt{f(x)})\\
%         &= \sum_{i=1}^n q_i(x) (\sqrt{f(x)}\sqrt{f(x^i)} - f(x))\\
%         &= \frac 12 \sum_{i=1}^n q_i(x) (f(x^i) - f(x)) - \frac 12 \sum_{i=1}^n q_i(x) (\sqrt{f(x^i)} - \sqrt{f(x)})^2\\
%         &= \frac 12 L f(x) - \Gamma(\sqrt{f})(x). \qedhere
%     \end{align*}
% \end{proof}

% {\color{red} Xinayuan: Simpler proof:
\begin{proof}
    By the definition of carr\'e du champ operator, 
    \begin{align*}
        \Gamma(\sqrt{f}) = \frac{1}{2} \tp{L(\sqrt{f}^2) - 2\sqrt{f} L \sqrt{f}} = \frac{1}{2} Lf - \sqrt{f} L \sqrt{f}.
    \end{align*}
    This completes the proof.
\end{proof}
% }










\subsection{Influence and surface area of Boolean functions}

\begin{definition}
    Let $f: \Omega \to \mathbb{R}$ be a function.

    \begin{enumerate}[(1)]
        \item The \textit{total influence} of $f$ is defined as
        \begin{align*}
            I_w(f) := \E[\mu]{\Gamma_{w}(f)} = \E[\mu]{\norm{\nabla_w f}^2_2}.
        \end{align*}
        Also, let $I_{w,i}(f) := \E[\mu]{\Gamma_{w,i}(f)}$ denote the influence of the $i$-th coordinate.

        \item The \textit{(Boolean) surface area} of $f$ is defined as
        \begin{align*}
            A_w(f) := \E[\mu]{\sqrt{\Gamma_wf}} = \E[\mu]{\norm{\nabla_w f}_2}.
        \end{align*}

        \item The \textit{total squared influence} of $f$ is defined as
        \begin{align*}
            J_{w}(f) := \sum_{i=1}^n \tp{\E[\mu]{\abs{\delta_{w,i} f}}}^2 = \norm{\E[\mu]{|\nabla_w| f}}^2_2.
        \end{align*}
    \end{enumerate}
    
\end{definition}



\begin{example}[Uniform distribution]
    Consider the case where $\mu$ is uniform over $\Omega$. Let $q_i(x) = \frac{1}{2}$ and $w_i(x) = 1$ for all $i \in [n]$ and $x \in \Omega$.
    Let $f: \Omega \to \{\pm 1\}$ be a Boolean-valued function. 
    Below, we omit $w$ in the subscript since everything is unweighted.

    \begin{enumerate}[(1)]
        \item Total influence: We have
        \begin{align*}
            I_i(f) = \Pr[\mu]{f(x^i) \neq f(x)}, \quad I(f) = \sum_{i=1}^n I_i(f).
        \end{align*}

        \item Surface area: We have
        \begin{align*}
            A(f) = \E[\mu]{\sqrt{s(f)}}
        \end{align*}
        where $s(f)$ denotes the sensitivity of $f$; namely, $s(f)(x) = |\{i \in [n]: f(x^i) \neq f(x)\}|$ for each $x \in \Omega$.
    
        \item Total squared influence: We have
        \begin{align*}
            J(f) = \sum_{i=1}^n \Pr[\mu]{f(x^i) \neq f(x)}^2 = \sum_{i=1}^n I_i(f)^2.
        \end{align*}
    \end{enumerate}
\end{example}

\ZC{Let's keep it like this for now. I may use different notations in the intro for:
\begin{align*}
    \mathsf{Inf_i}(f) &:= \Pr[\mu]{f(x^i) \neq f(x)} \\
    \mathsf{Inf}(f) &:= \sum_{i=1}^n \Pr[\mu]{f(x^i) \neq f(x)} \\
    \mathsf{SurfArea}(f) &:= \E[\mu]{\sqrt{s(f)}} \\
    \mathsf{SqInf}(f) &:= \sum_{i=1}^n \Pr[\mu]{f(x^i) \neq f(x)}^2.
\end{align*}
}

\begin{example}[Canonical transition rates and weights]
    Let $\mu$ be an arbitrary distribution supported on $\Omega$. 
    Consider the canonical transition rates and weights defined as: for any $i \in [n]$ and $x \in \Omega$,
    \begin{align*}
        q^\star_i(x) := \frac 12\sqrt{\frac{\mu(x^i)}{\mu(x)}}
        \qquad\text{and}\qquad
        w^\star_i(x) := 2q^\star_i(x) = \sqrt{\frac{\mu(x^i)}{\mu(x)}}.
    \end{align*}
    Let $f: \Omega \to \{\pm 1\}$ be a Boolean-valued function. 
    Below, we use the superscript $\star$ to represent the canonical weights $w^\star$.

    \begin{enumerate}[(1)]
        \item Total influence: We have
        \begin{align*}
            I^\star_i(f) = \Pr[\mu]{f(x^i) \neq f(x)}, \quad I^\star(f) = \sum_{i=1}^n I^\star_i(f).
        \end{align*}
        The above equality is derived from definitions:
        \begin{align*}
            I^\star_i(f) &= \E[\mu]{q_i^\star(x)^2 (f(x^i) - f(x))^2}
            = \E[\mu]{\ind\{f(x^i)\neq f(x)\}\frac{\mu(x^i)}{\mu(x)}}
            \\&= \sum_{x^i\in \Omega} \ind\{f(x^i)\neq f(x^{ii})\}\mu(x^i)
            = \Pr[\mu]{f(x^i) \neq f(x)},
        \end{align*}
        where the third equality follows from the fact that $x^{ii}=x$ holds for $i\in[n]$ and $x\in \Omega$.
        \ZC{Add a short calculation to explain.}

        \item Surface area: We have
        \begin{align*}
            A^\star(f) = \E[\mu]{\sqrt{s_\mu(f)}},
        \end{align*}
        where $s_\mu(f)$ denotes a weighted sensitivity of $f$ defined by
        \begin{align*}
            s_\mu(f)(x) = \sum_{i=1}^n \frac{\mu(x^i)}{\mu(x)} \mathbf{1}\{f(x^i) \neq f(x)\}.
        \end{align*}
        If $b \le \sqrt{\mu(x^i)/\mu(x)} \le 1/b$ for all $x \in \Omega$ for some $b \in (0,1]$, then 
        \begin{align*}
            bA^\star(f) \le \E[\mu]{\sqrt{s(f)}} \le \frac 1b A^\star(f).
        \end{align*}
        \ZC{Add constants above.}

        \item Total squared influence: We have
        \begin{align*}
            J^\star(f) = \sum_{i=1}^n \E[\mu]{\sqrt{\frac{\mu(x^i)}{\mu(x)}} \mathbf{1}\{f(x^i) \neq f(x)\}}^2
        \end{align*}
        If $b \le \sqrt{\mu(x^i)/\mu(x)} \le 1/b$ for all $x \in \Omega$ for some $b \in (0,1]$, then 
        \begin{align*}
            b^2J^\star(f) \le \sum_{i=1}^n \Pr[\mu]{f(x^i) \neq f(x)}^2 \le \frac 1{b^2}J^\star(f).
        \end{align*}
        \ZC{Add constants above.}
    \end{enumerate}
    
\end{example}

By H\"older’s and Jensen’s inequalities, we obtain the following relations.
\begin{fact}
    For any function $f : \Omega \to \R$, we have that $J_w(f)\leq A_w(f)^2 \leq I_w(f)$.
\end{fact}

\subsection{Functional inequalities}
\begin{definition}
    The \textit{variance} of function $f: \Omega \to \mathbb{R}$ is
    \begin{align*}
        \Var[\mu]{f} := \E[\mu]{f^2} - (\E[\mu]{f})^2.
    \end{align*}
    The \textit{entropy} of function $f: \Omega \to \mathbb{R}_+$ is
    \begin{align*}
        \Ent[\mu]{f} := \E[\mu]{f \log f} - \E[\mu]{f} \log \E[\mu]{f}.
    \end{align*}
\end{definition}

\begin{definition}[Poincar\'e inequality]
    We say that the Poincar\'e inequality holds with constant $\lambda > 0$ if for any $f: \Omega \to \R$,
    \begin{align}\label{eq:Poincare}
        \Var[\mu]{f} \leq \frac 1\lambda \E[\mu]{\Gamma (f)}, \tag{Poincar\'e}
    \end{align}
    where $\Gamma(f)$ is the unweighted carr\'e du champ operator, i.e. $w_i \equiv 1$ for all $i\in [n]$.
    Or equivalently, $\Var[\mu]{P_t f}\leq e^{-2\lambda t}\Var[\mu]{f}$ holds.
\end{definition}

\begin{definition}[Log-Sobolev inequality]
    We say that the log-Sobolev inequality holds with constant $\rho_{\LSI} > 0$ if for any $f: \Omega \to \R$,
    \begin{align}\label{eq:LSI}
        \Ent[\mu]{f^2} \leq \frac 2{\rho_{\LSI}} \E[\mu]{\Gamma (f)}. \tag{LSI}
    \end{align}
\end{definition}

\begin{definition}[Modified log-Sobolev inequality of gradient type]
    We say that the modified log-Sobolev inequality of gradient type holds with constant $\rho_{\MLSI} > 0$ if for any $f: \Omega \to \R_+$,
    \begin{align}\label{eq:mLSI-gradient}
        \Ent[\mu]{f} \leq \frac 2{\rho_{\MLSI}} \E[\mu]{\frac{\Gamma (f)}{f}}. \tag{mLSI-gradient}
    \end{align}
\end{definition}

\ZC{Use $\norm{\cdot}_{p,\mu}$}
\begin{theorem}[Hypercontractivity]\label{thm:hypercontractivity}
    If the log-Sobolev inequality holds with constant $\rho > 0$ and $1\leq p \leq q \leq \infty$ satisfy $\frac{q-1}{p-1}\leq \exp(2\rho t)$, then for all $f: \Omega \to \R$, we have $\norm{P_t f}_{q,\mu} \leq \norm{f}_{p,\mu}$.
\end{theorem}

\subsection{Gaussian isoperimetric profile}
\begin{definition}[Gaussian isoperimetric profile]\label{def:gaussian_isoperimetric_profile}
    We define the Gaussian isoperimetric profile $\+I: [0,1] \to \R_+$ as $\+I(x) = \varphi(\Psi^{-1}(x))$, where $\varphi$ and $\Phi$ are the Probability Density Function (PDF) and Cumulative Distribution Function (CDF) of the standard Gaussian distribution, respectively.
\end{definition}
\begin{fact}\label{fact:gaussian_isoperimetric_profile_ineq}
    For any $a,b\in [0,1]$, we have that
    $$
        \+I(a)(\+I(b)-\+I(a)-\+I'(a)(b-a))\geq -(b-a)^2 .
    $$
\end{fact}

\section{Bakry--\'Emery Conditions and Gradient Estimates}

\subsection{BE theory, weak and strong gradient estimates}
For $f,g: \Omega \to \R$, we write $f\leq g$ for brevity if $f(x)\leq g(x)$ for all $x\in \Omega$.
\begin{definition}[Bakry--\'Emery conditions]
    Let $\mu$ be a distribution supported on $\Omega = \{\pm 1\}^n$ and $\{q_i: \Omega \mapsto \R_+\}_{i\in[n]}$ be transition rates reversible for $\mu$.
    Suppose $\{w_i:\Omega \to \R_+\}_{i\in[n]}$ is a collection of weight functions.
    \begin{itemize}
        \item We say that the \textit{Bakry--\'Emery condition} holds with constant $\rho_{\BE} > 0$ if for any function $f: \Omega \to \mathbb{R}$, it holds
        \begin{align}\label{eq:BE}
            \Gamma_{2;w}(f) \geq \rho_{\BE} \cdot \Gamma_w(f); \tag{BE}
        \end{align}
        \item We say that the \textit{reinforced Bakry--\'Emery condition} holds with constant $\rho_{\RBE} > 0$ if for any function $f: \Omega \to \R$, it holds 
        \begin{align}\label{eq:rBE}
            \Gamma_{2;w}(f) - \Gamma(\sqrt{\Gamma_w(f)}) \ge \rho_{\RBE} \cdot \Gamma_w(f). \tag{rBE}
        \end{align}
    \end{itemize}
\end{definition}

The following comparison follows immediately from the non-negativity of $\Gamma$.
\begin{fact}
If~\ref{eq:rBE} holds with constant $\rho > 0$, then~\ref{eq:BE} holds with the same constant $\rho$.
\end{fact}




\begin{definition}[Weak and strong gradient estimate]
    Let $\mu$ be a distribution supported on $\Omega = \{\pm 1\}^n$ and $\{q_i: \Omega \mapsto \R_+\}_{i\in[n]}$ be transition rates reversible for $\mu$.
    Suppose $\{w_i:\Omega \to \R_+\}_{i\in[n]}$ is a collection of weight functions. 
    \begin{itemize}
        \item We say that the \emph{Weak Gradient Estimate (WGE)} holds with constant $\rho_{\WGE} > 0$ if for any function $f:\Omega \to \R$ and $t\geq 0$, 
        \begin{align}\label{eq:WGE}
            \Gamma_w(P_t f) \leq e^{-2\rho_{\WGE} t} P_t(\Gamma_w(f)) \quad\iff\quad \norm{\nabla_w P_t f}^2_2 \leq e^{-2\rho_{\WGE} t} P_t(\norm{\nabla_w f}^2_2); \tag{WGE}
        \end{align}
        \item We say that the \emph{Strong Gradient Estimate (SGE)} holds with constant $\rho_{\SGE} > 0$ if for any function $f: \Omega \to \R$ and $t\geq 0$, 
        \begin{align}\label{eq:SGE}
            \sqrt{\Gamma_w(P_t f)} \leq e^{-\rho_{\SGE} t} P_t  \sqrt{\Gamma_w(f)} \quad\iff\quad \norm{\nabla_w P_t f}_2 \leq e^{-\rho_{\SGE} t} P_t(\norm{\nabla_w f}_2); \tag{SGE}
        \end{align}
    \end{itemize}
\end{definition}

\begin{fact}\label{fact:SGE-to-WGE}
If~\ref{eq:SGE} holds with constant $\rho > 0$, then~\ref{eq:WGE} holds with the same constant $\rho$.
\end{fact}

\begin{proof}
    By~\ref{eq:SGE}, it holds that for any $f : \Omega \to \mathbb{R}$ and $t \ge 0$,
    \begin{align*}
         \norm{\nabla_w P_t f}^2_2  \le \e^{-2\rho t} \tp{P_t \norm{\nabla_w f}_2}^2 \le \e^{-2\rho t}P_t \tp{\norm{\nabla_w f}^2_2},
    \end{align*}
    where the last inequality follows from the Cauchy--Schwarz inequality: $\tp{P_t g}^2 \le P_t (g^2)$ for any function $g : \Omega \to \mathbb{R}$.
\end{proof}


\begin{proposition}\label{prop:equivalence_BE_GE}
    The following equivalences hold.
    \begin{itemize}
        \item The Bakry--\'Emery condition with constant $\rho$ iff the weak gradient estimate with constant $\rho$.
        \item The reinforced Bakry--\'Emery condition with constant $\rho$ iff the strong gradient estimate with constant $\rho$.
    \end{itemize}
\end{proposition}

\begin{lemma}\label{lem:gamma2-formula}
    For any function $f: \Omega \to \R$, it holds
    \begin{align*}
        \Gamma_{2;w}(f) - \Gamma(\sqrt{\Gamma_w(f)}) = \sqrt{\Gamma_w(f)} L \sqrt{\Gamma_w(f)} - \Gamma_w(f,Lf).
    \end{align*}
\end{lemma}

\begin{proof}
    % \ZC{Use \Cref{lem:chain-rule}}
    Straightforward calculation shows that
    \begin{align*}
        \Gamma_{2;w}(f) - \Gamma(\sqrt{\Gamma_w(f)}) &= \frac 12 L \Gamma_w(f) - \Gamma_w(f,Lf) - \Gamma(\sqrt{\Gamma_w(f)}) &&\text{(by definition of $\Gamma_{2;w}$)}
        \\&= \sqrt{\Gamma_w(f)} L \sqrt{\Gamma_w(f)} - \Gamma_w(f,Lf). &&\text{(by \Cref{lem:chain-rule})} \qedhere
    \end{align*}
\end{proof}
Now we prove \Cref{prop:equivalence_BE_GE} based on \Cref{lem:gamma2-formula}.


\begin{proof}[Proof of \Cref{prop:equivalence_BE_GE}]
\ZC{GPT: The converse directions require differentiating WGE and SGE at \(t=0\).
Notation change: $F(s) \to F_s$ and $G(s) \to G_s$, because both are operators with argument $f$.
}
    Fix $f : \Omega \to \R$ and $t \ge 0$. Let $F(s) = P_s \Gamma_w(P_{t-s} f)$. The derivative $\frac{\dif }{\dif s} F(s)$ satisfies
    \begin{align*}
        \frac{\dif }{\dif s} F(s) = L P_s \Gamma_w(P_{t-s} f) + P_s \frac{\dif}{\dif s} \Gamma_w(P_{t-s} f) &= L P_s \Gamma_w(P_{t-s} f) -2 P_s \Gamma_w(P_{t-s}f, LP_{t-s} f)\\
        &= 2 P_s \Gamma_{2;w}(P_{t-s} f).
    \end{align*}
    This proves the first part. Similarly, let $G(s) = P_s \sqrt{\Gamma_w(P_{t-s} f)}$. At points where $\Gamma_w(P_{t-s}f)=0$, the following calculation is justified by replacing $\sqrt{\Gamma_w}$ with $\sqrt{\Gamma_w+\varepsilon}$ and letting $\varepsilon\downarrow0$. The derivative satisfies
    \begin{align}\label{eq:dG-1}
    \frac{\dif}{\dif s} G(s) = P_s \tp{L \sqrt{\Gamma_w(P_{t-s} f)} - \frac{\Gamma_w(P_{t-s} f, L P_{t-s} f)}{\sqrt{\Gamma_w(P_{t-s} f)}}}.
    \end{align} 
    By the definition of the second-order carr\'e du champ operator $\Gamma_{2;w}$, it holds
    \begin{align}\label{eq:dG-2}
        \Gamma_w(P_{t-s} f, L P_{t-s} f) = \frac{1}{2} L \Gamma_w(P_{t-s} f) - \Gamma_{2;w}(P_{t-s} f). 
    \end{align}
    By~\Cref{lem:chain-rule}, the quantity $\sqrt{\Gamma_w(P_{t-s} f)} L\sqrt{\Gamma_w(P_{t-s} f)}$ can be written as
    \begin{align}\label{eq:dG-3}
        \sqrt{\Gamma_w(P_{t-s} f)} L\sqrt{\Gamma_w(P_{t-s} f)} = \frac{1}{2} L \Gamma_w(P_{t-s} f) - \Gamma(\sqrt{\Gamma_w(P_{t-s} f)}).
    \end{align}
    % \ZC{Cite the chain rule for \cref{eq:dG-3}? (\Cref{lem:gamma2-formula})}
    Combining~\ref{eq:dG-1},~\ref{eq:dG-2} and~\ref{eq:dG-3}, we have
    \begin{align*}
    \frac{\dif}{\dif s} G(s) = P_s\tp{\frac{\Gamma_{2;w}(P_{t-s} f) - \Gamma(\sqrt{\Gamma_w(P_{t-s} f)})}{\sqrt{\Gamma_w(P_{t-s} f)}}}.
    \end{align*}
    This proves the second part.
\end{proof}

\begin{remark}[Comparison with continuous space]
    Let $w \equiv 1$.
    For continuous spaces like $\mathbb{R}^n$, the Bakry--\'Emery condition is also known as the curvature-dimension $CD(\rho,\infty)$, which is also equivalent to the reinforced Bakry--\'Emery condition. However, in the discrete setting, the reinforced version is strictly stronger. \ZC{Need to expand.}
\end{remark}

\begin{proposition}[Contraction in total influence and surface area]
    %Let $\rho > 0$ be a constant.
    \quad
    \begin{itemize}
        \item If~\ref{eq:WGE} holds with constant $\rho > 0$, then the total influence contracts along the heat semigroup with rate $2\rho$; that is, for any function $f: \Omega \to \R$ and $t \ge 0$,
        \begin{align*}
            I_w(P_t f) \leq e^{-2\rho t} I_w(f); \tag{$I$-contraction}
        \end{align*}
        
        \item If~\ref{eq:SGE} holds with constant $\rho > 0$, then the surface area contracts along the heat semigroup with rate $\rho$; that is, for any function $f: \Omega \to \R$ and $t \ge 0$,
        \begin{align*}
            A_w(P_t f) \leq e^{-\rho t} A_w(f). \tag{$A$-contraction}
        \end{align*}   
    \end{itemize}
\end{proposition}

\begin{proof}
By the definition of the total influence, it holds
\begin{align*}
    I_w(P_t f) = \E[\mu]{\Gamma_w(P_t f)} \le \e^{-2 \rho t} \E[\mu]{P_t \Gamma_w(f)} = \e^{-2 \rho t} \E[\mu]{\Gamma_w(f)}.
\end{align*}
This proves the contraction of the total influence. 
Similarly, by the definition of the surface area, 
\begin{align*}
    A_w(P_t f) = \E[\mu]{\sqrt{\Gamma(P_t f)}} \le \e^{-\rho t} \E[\mu]{P_t \sqrt{\Gamma(f)}} = \e^{-\rho t} A_w(f).
\end{align*}
This proves the contraction of the surface area.
\end{proof}

\subsection{Coordinate-wise strong gradient estimate}
Suppose $v: \Omega \to \R^n$ is a vector field; for example, given any function $f: \Omega \to \R$, the discrete gradient $\nabla_w f$ and the absolute discrete gradient $\abs{\nabla_w}f$ are both vector fields.
Suppose $T$ is an operator mapping a function $f: \Omega \to \R$ to another function $Tf: \Omega \to \R$.
With a slight abuse of notation, we extend $T$ to apply to vector fields entrywisely; specifically, $T v: \Omega \to \R^n$ is again a vector field defined as
\begin{align*}
    T v := (T v_1,\dots,T v_n)
    % \begin{pmatrix}
    %     T v_1 \\
    %     \vdots \\
    %     T v_n
    % \end{pmatrix}
\end{align*}
where $v_1,\dots,v_n$ are components of $v$.

\begin{definition}[Coordinate-wise strong gradient estimate]
    Let $\mu$ be a distribution supported on $\Omega = \{\pm 1\}^n$ and $\{q_i: \Omega \mapsto \R_+\}_{i\in[n]}$ be transition rates reversible for $\mu$.
    Suppose $\{w_i:\Omega \to \R_+\}_{i\in[n]}$ is a collection of weight functions. 
    
    We say that the \emph{Coordinate-wise Strong Gradient Estimate (CSGE)} holds with constant $\rho_{\CSGE} > 0$ if for any $\tau \in [0,\infty]$, any function $f: \Omega \to \R$, and $t\geq 0$, 
        \begin{align}\label{eq:CSGE}
            \norm{P_\tau \abs{\nabla_w}P_t f}_2 \leq e^{-\rho_{\CSGE} t}\norm{P_{\tau+t} \abs{\nabla_w}f}_2. \tag{CSGE}
        \end{align}
        We record two particularly important subcases:
        \begin{itemize}
            \item ($\tau = 0$) For any function $f: \Omega \to \R$ and $t\geq 0$,  
            \begin{align}\label{eq:0-CSGE}
                \norm{\nabla_w P_t f}_2 \leq e^{-\rho_{\CSGE} t}\norm{P_t \abs{\nabla_w}f}_2. \tag{$0$-CSGE}
            \end{align}

            \item ($\tau = \infty$) For any function $f: \Omega \to \R$ and $t\geq 0$,  
            \begin{align}\label{eq:infty-CSGE}
                \norm{\E[\mu]{\abs{\nabla_w}P_t f}}_2 \leq e^{-\rho_{\CSGE} t}\norm{\E[\mu]{\abs{\nabla_w}f}}_2. \tag{$\infty$-CSGE}
            \end{align}
        \end{itemize}
\end{definition}

\begin{fact}
If~\ref{eq:0-CSGE} holds with constant $\rho > 0$, then both~\ref{eq:WGE} and~\ref{eq:SGE} hold with the same constant $\rho$.
\end{fact}

\begin{proof}
    By~\eqref{eq:0-CSGE} and Minkowski inequality,
    \begin{align*}
        \norm{\nabla_w P_t f}_2 \le \e^{-\rho t} \norm{P_t \abs{\nabla_w} f}_2 \le \e^{-\rho t} P_t \norm{\nabla_w f}_2.
    \end{align*}
    Therefore,~\eqref{eq:SGE} holds with the constant $\rho$. 
    By~\Cref{fact:SGE-to-WGE},~\eqref{eq:WGE} also holds with the same constant $\rho$.
\end{proof}


\begin{proposition}[Contraction in total squared influence]
    %Let $\rho > 0$ be a constant.
    If~\ref{eq:infty-CSGE} holds with constant $\rho > 0$, then the total squared influence contracts along the heat semigroup with rate $2\rho$; that is, for any function $f: \Omega \to \R$ and $t \ge 0$,
    \begin{align*}
        J_w(P_t f) \leq e^{-2\rho t} J_w(f). \tag{$J$-contraction}
    \end{align*} 
\end{proposition}

\begin{proof}
    Assume~\ref{eq:infty-CSGE} holds with constant $\rho$. For any $f:\Omega \to \mathbb{R}$ and $t \ge 0$, it holds
    \begin{align*}
        J_w(P_t f) = \norm{\E[\mu]{\abs{\nabla_w} P_t f}}_2^2 \le \e^{-2\rho t} \norm{\E[\mu]{\abs{\nabla_w} f}}_2^2 = \e^{-2 \rho t} J_w(f).
    \end{align*}
    The last two items can be proved similarly.
\end{proof}


\subsection{Functional inequalities from gradient estimates}
\begin{definition}
    For $b \in (0,1]$, we say the weight functions $\{w_i:\Omega \to \R_+\}_{i\in[n]}$ are $b$-bounded if for all $i \in [n]$ and $x \in \Omega$, it holds
    \begin{align}\label{eq:w_lower_bound}
        b \le w_i(x) \le \frac{1}{b}.
    \end{align}
\end{definition}

\begin{proposition}[Gradient estimate and functional inequalities]
    Let $\mu$ be a distribution supported on $\Omega = \{\pm 1\}^n$ and $\{q_i: \Omega \mapsto \R_+\}_{i\in[n]}$ be transition rates reversible for $\mu$.
    Suppose $\{w_i:\Omega \to \R_+\}_{i\in[n]}$ is a collection of weight functions 
    that are $b$-bounded for some $b \in (0,1]$.
    \begin{itemize}
        \item If~\ref{eq:SGE} holds with $\rho > 0$, then~\ref{eq:mLSI-gradient} holds with constant $4b^2\rho$;
        \item If~\ref{eq:WGE} holds with $\rho > 0$, then~\ref{eq:Poincare} holds with constant $b^2\rho$.
    \end{itemize}
\end{proposition}

\begin{proof}
    By a standard calculation,
    \begin{align*}
    \Ent[\mu]{f} = \int_0^{+\infty} \E[\mu]{\Gamma(P_t f, \log P_t f)} \dif t  \overset{(\star)}{\le}  \int_0^{+\infty} \E[\mu]{\frac{\Gamma(P_t f)}{P_t f}} \dif t \le \frac{1}{b} \cdot \int_0^{+\infty} \E[\mu]{\frac{\Gamma_w(P_t f)}{P_t f}} \dif t,
    \end{align*}
    where the inequality $(\star)$ follows from $(s-t)\tp{\log s -\log t} \le \frac{(s-t)^2}{2}\tp{\frac{1}{s} + \frac{1}{t}}$. By~\ref{eq:SGE},
    \begin{align*}
        \Gamma_w(P_t f) \le \e^{-2\rho t} \tp{P_t \sqrt{\Gamma_w(f)}}^2 \le \e^{-2\rho t} \tp{P_t \tp{\frac{\Gamma_w (f)}{f}}} \tp{P_t f},
    \end{align*}
    where the last inequality follows from Cauchy-Schwarz inequality. Hence,
    \begin{align*}
    \Ent[\mu]{f} \le \frac{1}{b} \int_0^{+\infty} \e^{-2 \rho t} \E[\mu]{\frac{\Gamma_w (f)}{f}} \dif t \le \frac{1}{2 b^2 \rho} \cdot \E[\mu]{\frac{\Gamma(f)}{f}}.
    \end{align*}
    This proves the first part. For the second part, we use a similar proof strategy. Specifically,
    \begin{align*}
    \Var[\mu]{f} = 2\int_0^{+\infty} \E[\mu]{\Gamma(P_t f)} \dif t &\le \frac{2}{b} \cdot \int_0^{+\infty} \E[\mu]{\Gamma_w(P_t f)} \dif t \\
    (\text{By~\ref{eq:WGE}})\quad & \le \frac{2}{b} \cdot \int_0^{+\infty} \e^{-2\rho t} \E[\mu]{\Gamma_w(f)} \dif t \le \frac{1}{b^2\rho} \cdot \E[\mu]{\Gamma(f)}.
    \end{align*}
    This establishes the desired Poincar\'e inequality.
\end{proof}

\ZC{Add a figure for implications among: BE, rBE, WGE, SGE, CGE, I/A/J-contraction, PI, mLSI}

\ZC{More implications}